% !TEX program = pdflatex
\documentclass{bmstu}

% --- Пакеты ---
\usepackage{xcolor}
\usepackage{tabularx}
\usepackage{longtable}
\usepackage{booktabs}
\usepackage{lscape}
\usepackage{afterpage}
\usepackage{amsmath}
\usepackage{amssymb}
\usepackage{csquotes}
\usepackage{xifthen}
\usepackage{graphicx}
\usepackage{float}
\usepackage{wrapfig}
\usepackage{tikzscale}
\usepackage{tikz}
\usetikzlibrary{positioning, shapes.geometric, arrows.meta, decorations.pathreplacing}
\usepackage{pgfplots}
\pgfplotsset{compat=newest}
\usepackage{enumitem}
\usepackage[titletoc]{appendix}
\usepackage{stackengine}
\usepackage[normalem]{ulem}
\usepackage{multirow}


% --- Настройка списков ---
\setlist[itemize]{label=---, wide, labelindent=\parindent, labelsep=0.5em, itemsep=2pt, parsep=0pt, topsep=4pt}
\makeatletter
\@ifundefined{AddEnumerateCounter}{}{\AddEnumerateCounter{\asbuk}{\russian@alph}{щ}}
\makeatother
\setlist[enumerate,1]{label=\arabic*), wide, labelindent=\parindent, labelsep=0.5em, itemsep=2pt, parsep=0pt, topsep=4pt}
\setlist[enumerate,2]{label=\asbuk*), wide, labelindent=\parindent, labelsep=0.5em, itemsep=2pt, parsep=0pt, topsep=4pt}
\newcommand{\blank}[1][1.5cm]{\underline{\hspace{#1}}}

\begin{document}

% ======================================================================
%  ТИТУЛЬНЫЙ ЛИСТ
% ======================================================================
\makeresearchtitle
{Информатика и системы управления (ИУ)} % факультет
{Информационные системы и телекоммуникации (ИУ3)} % кафедра
{Проектирование структурированной кабельной системы офиса лаборатории тестирования устройств IoT} % тема
{***} % студент / группа
{***} % руководитель
{} % консультант

% ======================================================================
%  РЕФЕРАТ
% ======================================================================
\begin{essay}{структурированная кабельная система, иерархическая звезда, проектирование СКС, телекоммуникационный шкаф, видеонаблюдение, активное сетевое оборудование, маршрутизация, импортозамещение}

Объектом разработки является структурированная кабельная система (СКС) этажа лаборатории тестирования устройств IoT.

Цель работы~--- разработка СКС этажа по архитектуре иерархической звезды с одной серверной и одной кроссовой комнатой.

В~процессе работы были выполнены расчеты количества информационных розеток, емкости коммутационных панелей и плотности заполнения телекоммуникационных шкафов. Проведен расчет трафика камер видеонаблюдения и нагрузки на медиасервер и операторскую в~штатном (MCU) и аварийном (SFU) режимах. Обоснована пропускная способность каналов связи и осуществлен выбор камер с~поддержкой кодеков H.265/H.265+.

Спроектирована физическая инфраструктура сети: рассчитаны контрольные длины кабельных линий, сечения и уровни заполнения кабельных каналов и лотков, а~также механические нагрузки на элементы крепления. Выполнен тепловой расчет, включающий оценку тепловыделения, параметров охлаждения и требуемой емкости ИБП для шкафов серверной и кроссовой. Разработана и обоснована компоновка телекоммуникационного оборудования с~учетом требований к~заземлению.

В~результате выполнения работы разработано техническое задание на СКС в~соответствии с~требованиями ГОСТ~34.602-2020 и ГОСТ~Р~15.301-2016. Осуществлен подбор активного сетевого оборудования из Единого реестра российской радиоэлектронной продукции Минпромторга~РФ и~произведен проверочный расчет спроектированной вычислительной сети. Подготовлен полный комплект отчетной документации, включающий расчетно-пояснительную записку, чертеж схемы сети объекта на~листе формата~А2 по~стандартам ЕСКД, схему логической топологии и ведомость покупных изделий.
\end{essay}

% ======================================================================
%  ОПРЕДЕЛЕНИЯ, ОБОЗНАЧЕНИЯ И СОКРАЩЕНИЯ
% ======================================================================
\begin{abbreviations}
\definition{АКБ}{аккумуляторная батарея}
\definition{БП}{блок питания}
\definition{БРП}{блок распределения питания}
\definition{ГОСТ~Р}{государственный стандарт России}
\definition{ЕСКД}{единая система конструкторской документации}
\definition{ИБП}{источник бесперебойного питания}
\definition{КНС}{кабеленесущие системы}
\definition{КС}{кабельная система}
\definition{КПД}{коэффициент полезного действия}
\definition{МГТУ~им.~Н.\,Э.~Баумана}{Федеральное государственное автономное образовательное учреждение высшего образования <<Московский государственный технический университет имени Н.\,Э.~Баумана (национальный исследовательский университет)>>}
\definition{НИР}{научно-исследовательская работа}
\definition{РПЗ}{расчетно-пояснительная записка}
\definition{РФ}{Российская Федерация}
\definition{СКС}{структурированная кабельная система}
\definition{СКУД}{система контроля и управления доступом}
\definition{ТЗ}{техническое задание}
\definition{ТО}{техническое обслуживание}
\definition{ТС}{технические средства кабельной системы}
\definition{ЭРП}{этажный распределительный пункт}
\definition{IoT}{Internet of Things (интернет вещей)}
\definition{MCU}{Multipoint Control Unit (центральный узел микширования потоков)}
\definition{MUTOA}{Multi-User Telecommunications Outlet Assembly (многопользовательская телекоммуникационная розетка)}
\definition{PoE}{Power over Ethernet (питание по кабелю Ethernet)}
\definition{RJ-45}{Registered Jack~45}
\definition{SFU}{Selective Forwarding Unit (узел селективной ретрансляции потоков)}
\definition{U}{Unit}
\end{abbreviations}

% ======================================================================
%  СОДЕРЖАНИЕ
% ======================================================================
\maketableofcontents

% ======================================================================
%  ВВЕДЕНИЕ
% ======================================================================
\chapter*{ВВЕДЕНИЕ}
\addcontentsline{toc}{chapter}{ВВЕДЕНИЕ}

В~условиях стремительного развития цифровых технологий и роста объемов передаваемых данных надежная телекоммуникационная инфраструктура является критически важным элементом любого современного предприятия. Особые требования к~производительности, отказоустойчивости и безопасности сетей предъявляются в~специализированных испытательных центрах и лабораториях, где осуществляется работа с~большими массивами тестовых данных и непрерывный мониторинг процессов.

Актуальность разработки структурированной кабельной системы (СКС) для лаборатории тестирования устройств IoT обусловлена необходимостью создания единой, масштабируемой и высокоскоростной физической среды передачи сигналов. Современная СКС должна не~только обеспечивать информационный обмен между лабораторными стендами, операторскими рабочими местами и серверами, но~и~поддерживать интеграцию систем физической безопасности: видеонаблюдения высокого разрешения и~систем контроля и~управления доступом (СКУД) с~подачей электропитания по~технологии PoE+/PoE++. Дополнительную актуальность проекту придает необходимость реализации принципов импортозамещения при~выборе активного сетевого оборудования.

Объектом исследования и разработки является телекоммуникационная инфраструктура этажа лаборатории тестирования устройств IoT.

Предметом разработки выступают методы проектирования физического уровня локальных вычислительных сетей, принципы расчета кабельных трасс, энергетического баланса и параметров микроклимата коммутационных узлов.

Целью данной научно-исследовательской работы является проектирование масштабируемой и надежной структурированной кабельной системы этажа лаборатории на~базе архитектуры иерархической звезды, отвечающей современным стандартам пропускной способности и отказоустойчивости.

Для достижения поставленной цели в~рамках работы сформулированы и решены следующие задачи:
\begin{itemize}
    \item проведен анализ архитектурного плана и исходных данных объекта, осуществлено оптимальное размещение лабораторного, операторского и телекоммуникационного оборудования;
    \item выполнен расчет необходимого количества информационных портов, емкости коммутационных панелей и уровня заполнения телекоммуникационных шкафов в~серверной и кроссовой;
    \item произведен расчет пропускной способности каналов связи с~учетом трафика камер видеонаблюдения в~штатном (MCU) и аварийном (SFU) режимах;
    \item спроектированы кабельные трассы: рассчитаны контрольные длины линий, заполняемость кабельных каналов и лотков, а~также механические нагрузки на~элементы крепления;
    \item выполнен комплексный расчет параметров микроклимата и энергоснабжения, включающий оценку тепловыделения, выбор систем охлаждения и расчет требуемой емкости источников бесперебойного питания (ИБП);
    \item осуществлен подбор активного сетевого оборудования из~Единого реестра российской радиоэлектронной продукции Минпромторга~РФ с~последующим проверочным расчетом сети;
    \item разработано техническое задание на~СКС и подготовлен полный комплект проектной и конструкторской документации.
\end{itemize}


% ======================================================================
%  РАСЧЁТ ПРОПУСКНОЙ СПОСОБНОСТИ КАНАЛОВ СВЯЗИ
% ======================================================================

\chapter{Обследование объекта и порты}

\section{Общее число портов СКС}

Общее число портов СКС определяется как сумма портов по всем категориям помещений:
\begin{equation}
    N_{\text{общ}} = N_{\text{лаб}} + N_{\text{опер}} + N_{\text{проч}},
    \label{eq:n_total}
\end{equation}
где:
\begin{itemize}
    \item $N_{\text{лаб}}$~--- порты в~лабораториях;
    \item $N_{\text{опер}}$~--- порты в~операторской;
    \item $N_{\text{проч}}$~--- порты в~прочих помещениях.
\end{itemize}

\textbf{Порты в~лабораториях.} Согласно варианту~V10, под лаборатории отведено 2~помещения (№\,4 и~№\,11), в~которых размещено 6~лабораторных стендов (по 3~стенда на~лабораторию). В~каждом стенде устанавливается многопользовательская розетка (MUTOA), ёмкость которой определяется числом тестируемых IoT-устройств, портом узконаправленной камеры контроля и резервными портами для контрольно-измерительного оборудования (КИП).

Для помещения №\,4 (Лаб.\,1) ёмкость MUTOA одного стенда составляет:
\begin{equation}
    N_{\text{MUTOA\_лаб1}} = N_{\text{IoT}} + N_{\text{КИП}} + N_{\text{кам}} = 18 + 6 + 1 = 25 \text{ портов}.
    \label{eq:mutoa_lab1}
\end{equation}

Для помещения №\,11 (Лаб.\,2) ёмкость MUTOA одного стенда составляет:
\begin{equation}
    N_{\text{MUTOA\_лаб2}} = N_{\text{IoT}} + N_{\text{КИП}} + N_{\text{кам}} = 8 + 6 + 1 = 15 \text{ портов}.
    \label{eq:mutoa_lab2}
\end{equation}

Суммарная ёмкость MUTOA лабораторных стендов для помещения №\,4:
\begin{equation}
    N_{\text{лаб1\_MUTOA}} = 3 \times 25 = 75 \text{ портов}.
\end{equation}

Суммарная ёмкость MUTOA лабораторных стендов для помещения №\,11:
\begin{equation}
    N_{\text{лаб2\_MUTOA}} = 3 \times 15 = 45 \text{ портов}.
\end{equation}

Дополнительно в~каждой из 2~лабораторий размещаются IP-камера общего вида и контроллер СКУД, подключаемые к~портам PoE+. Это добавляет:
\begin{equation}
    N_{\text{лаб\_доп}} = 2 \times (1_{\text{камера}} + 1_{\text{СКУД}}) = 4 \text{ порта}.
\end{equation}

Итого портов в~лабораториях:
\begin{equation}
    N_{\text{лаб}} = 75 + 2 + 45 + 2 = 124 \text{ порта}.
\end{equation}

(Проверка: помещение №\,4: $3 \times 25 + 1 + 1 = 77$ портов; помещение №\,11: $3 \times 15 + 1 + 1 = 47$ портов; $77 + 47 = 124$ порта.)

\textbf{Порты в~операторской.} Операторская размещается в~помещении №\,20 (<<Опенспейс>>) и включает 6~рабочих мест (постов). На~каждом посту размещается 1~моноблок (1~порт PoE++), 1~рабочая станция (1~порт) и 1~резервная розетка. Также в~общей зоне операторской устанавливаются 2~дополнительные розетки:
\begin{equation}
    N_{\text{опер}} = 6 \times (1_{\text{мон}} + 1_{\text{РС}} + 1_{\text{рез}}) + 2_{\text{общ}} = 18 + 2 = 20 \text{ портов}.
    \label{eq:n_oper}
\end{equation}

\textbf{Порты в~прочих помещениях.} Двумя розетками СКС оборудуются 2~помещения (№\,7, №\,24):
\begin{equation}
    N_{\text{проч}} = 2 \times 2 = 4 \text{ порта}.
    \label{eq:n_other}
\end{equation}

Итого общее число портов СКС:
\begin{equation}
    N_{\text{общ}} = 124 + 20 + 4 = 148 \text{ портов}.
    \label{eq:n_total_calc}
\end{equation}

\textbf{Разграничение по типам PoE:}
\begin{itemize}
    \item PoE++ (IEEE~802.3bt, до~90~Вт): моноблоки в~операторской~--- 6~портов;
    \item PoE+ (IEEE~802.3at, до~30~Вт): камеры контроля стендов (6~портов) + камеры общего вида (2~порта) + контроллеры СКУД (2~порта)~--- 10~портов;
    \item Без PoE: IoT-устройства, резервные порты MUTOA, рабочие станции и резервные розетки операторской, розетки прочих помещений~--- 132~порта.
\end{itemize}

Итого: PoE++~--- 6~портов, PoE+~--- 10~портов, без~PoE~--- 132~порта.

% ---- Расчёт нагрузки от видео трафика ----
\chapter{Трафик, архитектура, выбор активного оборудования}

\section{Логическая архитектура и топология сети}

Согласно техническому заданию и современным корпоративным стандартам, сеть предприятия строится по классической трехуровневой иерархической модели (Cisco Three-Tier Model), что обеспечивает масштабируемость, отказоустойчивость и безопасность.

\subsection{Уровни иерархии и аплинки}

\begin{enumerate}
    \item \textbf{Уровень доступа (Access):} Располагается в кроссовой комнате (пом. №\,16). Задача уровня — терминация горизонтальных кабельных линий (148 портов) и обеспечение конечных устройств доступом в сеть с требуемыми параметрами электропитания (PoE+, PoE++). 
    
    \item \textbf{Уровень ядра (Core):} Размещается в серверной комнате этажа (пом. №\,19). Соединяется с коммутаторами уровня доступа кроссовой комнаты (пом. №\,16) напрямую посредством магистральных оптических каналов 10G SFP+. Ядро осуществляет высокоскоростную маршрутизацию между VLAN, связь с серверами, главным шлюзом сети (маршрутизатором) и стояками других этажей (№\,9 и №\,11).
\end{enumerate}

\subsection{Логическая адресация и VLAN}

Для изоляции широковещательных доменов, обеспечения безопасности и управления качеством обслуживания (QoS), логическая структура сети сегментируется на виртуальные локальные сети (VLAN). Применяется следующая схема адресации (блок $10.10.x.x$):

\begin{itemize}
    \item \textbf{VLAN 10 (Management):} $10.10.10.0/24$ — управление сетевым оборудованием (коммутаторы, маршрутизатор, ИБП).
    \item \textbf{VLAN 20 (Users):} $10.10.20.0/24$ — рабочие станции операторов и административного персонала.
    \item \textbf{VLAN 30 (IoT \& Labs):} $10.10.30.0/23$ — расширенная подсеть для тестируемых устройств IoT на стендах лабораторий (вмещает до 510 хостов).
    \item \textbf{VLAN 40 (Video \& Security):} $10.10.40.0/24$ — IP-камеры видеонаблюдения, медиасервер и контроллеры СКУД. Изоляция данного трафика критична для обеспечения работы режима MCU/SFU.
    \item \textbf{VLAN 50 (Servers):} $10.10.50.0/24$ — сервер CI/CD, NAS хранилище, интерфейсы управления серверами.
\end{itemize}

Схема логической топологии приведена в Приложении В.


\section{Расчёт нагрузки от видео трафика}

Согласно техническому заданию, во время тестов в~лабораториях нагрузка формируется в~основном за~счёт видео трафика, к~которому добавляется незначительный по~объёму поток команд и телеметрических данных.

Формула для расчёта видеотрафика без применения кодека (несжатый поток):
\begin{equation}
    P' = \frac{W \cdot H \cdot FPS \cdot Bit}{10^{6}} \quad [\text{Мбит/с}],
    \label{eq:bitrate_raw}
\end{equation}
где:
\begin{itemize}
    \item $W$~--- ширина изображения в~пикселях;
    \item $H$~--- высота изображения в~пикселях;
    \item $FPS$~--- количество кадров в~секунду;
    \item $Bit$~--- количество бит для кодирования одного пикселя (при цветовой модели YCbCr~4:2:2 принимается $Bit = 16$).
\end{itemize}

Расчёт несжатых потоков для применяемых камер:
\begin{equation}
    P'_{\text{узк\_MCU}} = \frac{3840 \cdot 2160 \cdot 30 \cdot 16}{10^{6}} = 3981{,}3 \text{ Мбит/с},
    \label{eq:raw_narrow}
\end{equation}
\begin{equation}
    P'_{\text{общ\_MCU}} = \frac{3840 \cdot 2160 \cdot 30 \cdot 16}{10^{6}} = 3981{,}3 \text{ Мбит/с},
    \label{eq:raw_overview_mcu}
\end{equation}
\begin{equation}
    P'_{\text{общ\_SFU}} = \frac{640 \cdot 480 \cdot 25 \cdot 16}{10^{6}} = 122{,}88 \text{ Мбит/с}.
    \label{eq:raw_overview_sfu}
\end{equation}

Для оптимизации использования пропускной способности сети видеопотоки подвергаются аппаратному сжатию непосредственно на IP-камерах с помощью высокоэффективного кодека H.265 (HEVC) / H.265+. Данный стандарт межкадрового сжатия позволяет уменьшить объем передаваемых данных в сотни раз (коэффициент компрессии $K_c$ достигает 300--500) за счет кодирования только изменяющихся участков сцены (векторов движения) при сохранении исходного визуального качества.

Исходя из рекомендаций производителей систем видеонаблюдения и калькуляторов битрейта для кодека H.265 при настройке постоянного битрейта (CBR) и средней динамичности сцены в помещениях лабораторий, эмпирические значения пропускной способности составляют:
\begin{itemize}
    \item Для разрешения 4 Мп (30~Гц) коэффициент сжатия $K_c \approx 442$. Требуемая полоса: $3981{,}3 / 442 \approx 9$~Мбит/с.
    \item Для разрешения 640$\times$480 (25~Гц) коэффициент сжатия $K_c \approx 123$ (ниже из-за накладных расходов на сетевые заголовки при малом разрешении). Требуемая полоса: $122{,}88 / 123 \approx 1$~Мбит/с.
\end{itemize}

С~учетом этого, расчетные битрейты сжатых видеопотоков от одной камеры составляют:
\begin{itemize}
    \item $P_{\text{узк\_MCU}} = 9$~Мбит/с~--- нагрузка на узконаправленную камеру (4 Мп 30~Гц) в~режиме MCU;
    \item $P_{\text{общ\_MCU}} = 9$~Мбит/с~--- нагрузка на камеру общего вида (4 Мп 30~Гц) в~режиме MCU;
    \item $P_{\text{общ\_SFU}} = 1$~Мбит/с~--- нагрузка на камеру общего вида (640$\times$480 25~Гц) в~аварийном режиме SFU.
\end{itemize}

% ---- Расчёт нагрузки на коммутаторы доступа ----
\section{Расчёт нагрузки на коммутаторы доступа}

В~штатном режиме MCU на каждую лабораторию приходится нагрузка от 3~узконаправленных камер и 1~обзорной камеры:
\begin{equation}
    P_{\text{лаб1\_MCU}} = 3 \times P_{\text{узк\_MCU}} + 1 \times P_{\text{общ\_MCU}} = 3 \times 9 + 9 = 36 \text{ Мбит/с}.
    \label{eq:p_lab1}
\end{equation}
\begin{equation}
    P_{\text{лаб2\_MCU}} = 3 \times P_{\text{узк\_MCU}} + 1 \times P_{\text{общ\_MCU}} = 3 \times 9 + 9 = 36 \text{ Мбит/с}.
    \label{eq:p_lab2}
\end{equation}

Суммарный видеотрафик от 2~лабораторий, поступающий на коммутаторы доступа:
\begin{equation}
    P_{\text{доступа\_MCU}} = 36 + 36 = 72 \text{ Мбит/с}.
    \label{eq:p_access_mcu}
\end{equation}

В~аварийном режиме SFU направленные камеры отключаются, а~обзорные камеры переключаются в~режим 640$\times$480:
\begin{equation}
    P_{\text{доступа\_SFU}} = 2 \times P_{\text{общ\_SFU}} = 2 \times 1 = 2 \text{ Мбит/с}.
    \label{eq:p_access_sfu}
\end{equation}

Проверка: 72~Мбит/с $<$ 1~Гбит/с (производительность портов 1G обеспечена с~большим запасом).

% ---- Расчёт нагрузки на медиасервер ----
\section{Расчёт нагрузки на медиасервер}

\textbf{Штатный режим MCU:}
\begin{itemize}
    \item Входящий поток от лабораторий: $P_{\text{медиа\_вх}} = 72$~Мбит/с;
    \item Исходящий поток на 6~моноблоков (4 Мп 30~Гц): $P_{\text{медиа\_исх}} = 6 \times 9 = 54$~Мбит/с;
    \item Суммарная нагрузка на интерфейс медиасервера: $P_{\text{медиа\_MCU}} = 72 + 54 = 126$~Мбит/с.
\end{itemize}

\textbf{Аварийный режим SFU:}
\begin{itemize}
    \item Входящий поток от обзорных камер: $P_{\text{медиа\_вх}} = 2$~Мбит/с;
    \item Исходящий поток (трансляция 2~обзорных камер на 6~моноблоков): $P_{\text{медиа\_исх}} = 6 \times 2 \times 1 = 12$~Мбит/с;
    \item Суммарная нагрузка: $P_{\text{медиа\_SFU}} = 2 + 12 = 14$~Мбит/с.
\end{itemize}

% ---- Расчёт нагрузки на коммутатор ядра ----
\section{Расчёт нагрузки на коммутатор ядра}

Всего тестируемых IoT-устройств: $N_{\text{IoT}} = 78$~шт. ($3 \times 18 + 3 \times 8$). Значение скорости потока при прошивке одного устройства в 50~Мбит/с строится на решении обратной инженерной задачи, связывающей требования ТЗ и аппаратные ограничения IoT-микроконтроллеров. 

Средний объем бинарного образа встраиваемой ОС (ядро, файловая система, тестовые скрипты) составляет около 200~Мбайт (1600~Мбит). Согласно требованиям к регламентному времени, загрузка прошивки для стенда должна занимать от нескольких секунд до нескольких минут. При скорости 50~Мбит/с передача образа в 200~Мбайт займет 32 секунды, что идеально укладывается во временные рамки конвейера автоматизированной сборки. Кроме того, данная скорость обоснована физическими ограничениями интерфейсов плат: скорость поблочной записи встроенной Flash/eMMC памяти аппаратно ограничена 4--8~Мбайт/с (32--64~Мбит/с), поэтому передача на скоростях свыше 50~Мбит/с приведет лишь к переполнению буферов устройств и потере пакетов.

С учетом данного обоснования, пиковый трафик прошивок составляет:
\begin{equation}
    P_{\text{IoT\_тест}} = 78 \times 50 = 3900 \text{ Мбит/с}.
    \label{eq:p_iot}
\end{equation}

Оценка максимальной нагрузки на магистральный канал производится с учётом направления передачи данных (дуплексный режим).

Нагрузка при восходящем соединении (Uplink, от коммутаторов доступа к серверному сегменту) включает трансляцию видеопотоков на медиасервер ($P_{\text{доступа\_MCU}} = 72$~Мбит/с), передачу телеметрии от IoT-устройств ($P_{\text{IoT\_UP}} \approx 78 \times 1 = 78$~Мбит/с) и служебные запросы рабочих станций (принимаем с запасом $P_{\text{РМ\_вх}} = 100$~Мбит/с):
\begin{equation}
    P_{\text{uplink}} = P_{\text{доступа\_MCU}} + P_{\text{IoT\_UP}} + P_{\text{РМ\_вх}} = 72 + 78 + 100 = 250 \text{ Мбит/с}.
    \label{eq:p_uplink}
\end{equation}

Нагрузка при нисходящем соединении (Downlink, от сервера к абонентам) формируется пиковым трафиком параллельной прошивки IoT-плат ($P_{\text{IoT\_тест}} = 3900$~Мбит/с), трансляцией обработанного видео на рабочие места ($P_{\text{медиа\_исх}} = 54$~Мбит/с) и скачиванием документации с NAS-сервера (примем $P_{\text{NAS\_исх}} = 500$~Мбит/с):
\begin{equation}
    P_{\text{downlink}} = P_{\text{IoT\_тест}} + P_{\text{медиа\_исх}} + P_{\text{NAS\_исх}} = 3900 + 54 + 500 = 4454 \text{ Мбит/с}.
    \label{eq:p_downlink}
\end{equation}

Итоговая пиковая пропускная способность магистрального линка:
\begin{equation}
    P_{\text{магистраль}} = \max(250,\; 4454) = 4454 \text{ Мбит/с} \approx 4{,}45 \text{ Гбит/с}.
    \label{eq:p_core}
\end{equation}

Проверка: 4,45~Гбит/с $>$ 1~Гбит/с, следовательно, для магистральных каналов сервера и коммутаторов требуется использование портов 10G~SFP+ с поддержкой LACP.





% ---- Выбор коммутаторов доступа ----
\section{Выбор коммутаторов доступа}

Для подключения 148 абонентских портов СКС выбраны коммутаторы уровня доступа из Единого реестра российской радиоэлектронной продукции Минпромторга РФ. С целью оптимизации стоимости проекта и рационального использования монтажного пространства телекоммуникационного шкафа применена комбинированная схема распределения абонентов: устройства, не требующие дистанционного питания (всего 132 порта), частично подключаются к свободным интерфейсам коммутаторов с поддержкой PoE. В соответствии с требованиями к отказоустойчивости топологии, коммутаторы доступа объединяются в единый логический стек. Это позволяет управлять ими как одним виртуальным устройством и организовать агрегированные каналы связи (LACP) до ядра. 

Выбор оборудования, согласно технической документации производителей, разделен по требуемым стандартам питания:

\begin{enumerate}
    \item \textbf{Для устройств с питанием PoE++ (6 портов)}: выбран коммутатор Aquarius AQ-N3000-48P4Y2Q (48 портов 1G PoE++, 4 порта 25G SFP28, 2 порта 40G QSFP+, бюджет PoE 1080/2160 Вт), согласно паспорту устройства. Достаточно 1 коммутатора. Под моноблоки операторов задействуется 6 интерфейсов, а оставшиеся 42 порта выделяются для подключения оборудования без поддержки PoE и формирования технологического резерва.
    \item \textbf{Для оборудования с питанием PoE+ (10 портов)}: выбран коммутатор Eltex MES2448P (48 портов 1G PoE+, 4 порта 10G SFP+, бюджет PoE 720 Вт). Достаточно 1 коммутатора. Под видеокамеры и сетевые контроллеры СКУД задействуется 10 интерфейсов, свободные 38 портов маршрутизируются под стандартные абонентские розетки без PoE.
    \item \textbf{Для портов без PoE (основной массив)}: выбран коммутатор Eltex MES2448B (48 портов 1G RJ-45, 4 порта 10G SFP+), на основании официальных спецификаций. Из суммарно требуемых 132 портов без поддержки PoE, 80 портов уже обеспечены свободной емкостью вышеупомянутых коммутаторов Aquarius и Eltex MES2448P. Для покрытия оставшейся потребности (52 порта) требуется 2 коммутатора Eltex MES2448B, предоставляющих суммарно 96 портов.
\end{enumerate}

Для наглядности в таблице~\ref{tab:port_distribution} приведено итоговое распределение 148 абонентских портов по выбранным коммутаторам доступа, а также расчет технологического запаса (свободных портов).

\begin{table}[H]
    \centering
    \caption{Распределение портов по коммутаторам доступа}
    \label{tab:port_distribution}
    \renewcommand{\arraystretch}{1.2}
    \begin{tabularx}{\textwidth}{|>{\raggedright\arraybackslash}p{4cm}|>{\centering\arraybackslash}X|>{\centering\arraybackslash}X|>{\centering\arraybackslash}X|>{\raggedright\arraybackslash}p{3.2cm}|>{\centering\arraybackslash}p{2.2cm}|}
        \hline
        \textbf{Модель коммутатора} & \textbf{Всего портов} & \textbf{Занято (PoE)} & \textbf{Занято (без PoE)} & \textbf{Назначение (абоненты)} & \textbf{Свободно} \\
        \hline
        Aquarius AQ-N3000-48P4Y2Q & 48 & 6 (PoE++) & 42 & Моноблоки, абоненты б/п & 0 \\
        \hline
        Eltex MES2448P & 48 & 10 (PoE+) & 38 & Камеры, СКУД, абоненты б/п & 0 \\
        \hline
        Eltex MES2448B (№1) & 48 & 0 & 48 & Абоненты без PoE & 0 \\
        \hline
        Eltex MES2448B (№2) & 48 & 0 & 4 & Абоненты без PoE & 44 \\
        \hline
        \textbf{Итого} & \textbf{192} & \textbf{16} & \textbf{132} & \textbf{---} & \textbf{44} \\
        \hline
    \end{tabularx}
\end{table}

\section{Расчет аплинков уровня доступа}

В связи со стекированием коммутаторов уровня доступа, они работают как единое логическое устройство. Для исключения «узких мест» между стеком коммутаторов доступа и коммутатором ядра рассчитывается общая пропускная способность агрегированного восходящего канала (uplink). От каждого из 4 коммутаторов прокладывается по одной линии SFP+, формируя общую транковую группу (LACP). 

Суммарная емкость транка:
\begin{equation}
    P_{\text{транк}} = 4 \times 10\,000 = 40\,000 \text{ Мбит/с}.
\end{equation}

Пиковая генерация трафика в сети состоит из:
\begin{itemize}
    \item Трафик лабораторных стендов: $P_{\text{IoT\_тест}} = 3900$~Мбит/с.
    \item Трафик моноблоков операторов: $P_{\text{медиа\_исх}} = 54$~Мбит/с + $500$~Мбит/с (резерв).
    \item Трафик видеонаблюдения и СКУД: не более $100$~Мбит/с.
\end{itemize}

Суммарная пиковая нагрузка от уровня доступа к ядру:
\begin{equation}
    P_{\text{пик}} = 3900 + 554 + 100 = 4554 \text{ Мбит/с}.
\end{equation}

Оценка коэффициента утилизации агрегированного канала 40G:
\begin{equation}
    K_{\text{утил}} = \frac{4554}{40000} \approx 0{,}114 \text{ (запас 88{,}6\%)}.
\end{equation}

Таким образом, четырех линий 10G SFP+, объединенных в LACP-группу, более чем достаточно для обслуживания пиковых нагрузок стека доступа с огромным запасом под будущее масштабирование лаборатории.



% ---- Выбор коммутатора ядра ----
\section{Выбор коммутатора ядра}

Коммутатор ядра связывает коммутаторы доступа, серверное оборудование (медиасервер, NAS-хранилище, сервер сборки CI/CD) и~главный маршрутизатор.

Суммарная нагрузка на ядро составляет 4,45~Гбит/с, что~превышает лимит одного 1G~интерфейса. В~качестве коммутатора ядра выбран Eltex~MES5310-48 (24~порта 10G SFP+, 2~порта 100G QSFP28), полностью удовлетворяющий требованиям по~производительности, числу портов 10G и~импортозамещению.





\section{Выбор главного маршрутизатора}

Для обеспечения маршрутизации трафика локальной вычислительной сети и высокоскоростного подключения к сети связи общего пользования выбран сервисный маршрутизатор Eltex ESR-1200. В соответствии с техническим заданием устройство терминирует волоконно-оптические каналы связи со скоростью не менее 1 Гбит/с. Оборудование организует основной (к главной серверной здания) и резервный (к резервной серверной) вводы от слаботочного стояка СКС, что обеспечивает отказоустойчивость внешнего канала. Наличие интерфейсов 10G SFP+ позволяет создать агрегированный канал связи с коммутатором ядра. Маршрутизатор поддерживает протоколы динамической маршрутизации, межсетевое экранирование и трансляцию сетевых адресов (NAT), что полностью удовлетворяет проектным требованиям к изоляции и маршрутизации виртуальных подсетей (VLAN) лаборатории. 

\section{Выбор моноблоков}

Для оснащения рабочих мест операторов в помещении опенспейса (6 постов) выбрана энергоэффективная конфигурация моноблока «Гравитон» М75И (с процессором TDP 35 Вт и без дискретной графики). Пиковое энергопотребление такой конфигурации не превышает 65 Вт. Поскольку прямое питание моноблока по витой паре невозможно из-за отсутствия встроенного PoE-адаптера, для его подключения применяется PoE-сплиттер Planet POE-173S. Приобретение данного сплиттера, не входящего в Единый реестр, законно обосновывается в соответствии с 44-ФЗ и Постановлением Правительства РФ № 878 как закупка оборудования с отсутствующими российскими аналогами, удовлетворяющими жесткому техническому требованию поддержки стандарта IEEE 802.3bt Type 4 (мощность до 90 Вт на порт). Сплиттер принимает питание от коммутатора по стандарту IEEE 802.3bt Type 4 (PoE++) и преобразует его в 19 В постоянного тока, что позволяет использовать стандартный разъем питания моноблока и строго укладывается в гарантированный лимит 71,3 Вт на стороне питаемого устройства.


\section{Выбора видеокамер}

Система видеонаблюдения лабораторий строится на базе IP-камер с поддержкой дистанционного питания. Для детализированного контроля лабораторных стендов (6 шт.) выбраны узконаправленные IP-камеры QTECH QVC-MiR404Z532 с потребляемой мощностью 8 Вт. Для обеспечения общего контроля помещений лабораторий (2 шт.) применяются обзорные IP-камеры QTECH QVC-MiR803M с энергопотреблением 7 Вт. Обе модели камер поддерживают стандартизированное питание по технологии PoE (IEEE 802.3af/at). 



\section{Выбор серверного оборудования}

Вычислительная инфраструктура серверной комнаты базируется на высокопроизводительных стоечных решениях, монтируемых в 19-дюймовый телекоммуникационный шкаф. Оборудование подобрано с учетом обеспечения расчетной пропускной способности интерфейсов и резервирования аппаратных компонентов:

\begin{itemize}
    \item \textbf{Медиасервер:} Для централизованной обработки видеотрафика выбран сервер Гравитон С2242И. Устройство обладает достаточными вычислительными ресурсами для приема, перекодирования и трансляции индивидуальных видеопотоков высокого разрешения (4 Мп) в штатном режиме (MCU), а также для работы в аварийном режиме (SFU). Сетевая подсистема сервера включает оптические интерфейсы 10G SFP+, исключающие образование узких мест при приеме 72 Мбит/с и отдаче 54 Мбит/с видеотрафика в сторону коммутатора ядра. 
    \item \textbf{Сервер автоматизированной сборки (CI/CD):} Для компиляции тестового программного обеспечения и его одновременной загрузки на тестируемые устройства выбран сервер Гравитон С2122АА. Вычислительная платформа оснащена сетевыми адаптерами 10G SFP+, что гарантирует стабильную передачу пикового трафика прошивок (до 3900 Мбит/с для 78 IoT-устройств) без потерь и задержек. 
    \item \textbf{Сетевое хранилище (NAS):} Для кэширования дистрибутивов, хранения компонентов сборки и ведения логов тестирования применяется система хранения данных Qtech QSRV-260202. Накопитель поддерживает организацию аппаратного RAID-массива для повышения сохранности данных. Подключение к сети реализовано через оптический порт 10G SFP+ (или агрегированную группу портов) для повышения суммарной пропускной способности при обмене данными с сервером сборки и рабочими станциями.
\end{itemize}




\section{Анализ PoE-бюджетов коммутаторов}

Расчет энергетического баланса PoE производится в соответствии с проектными нормативами. Для обеспечения высокой надежности и исключения отключения портов при пусковых токах, расчет бюджета строится на основе максимальной мощности класса PoE (Class-based Power Allocation), к которому принадлежит подключаемое устройство. В бюджет коммутатора закладывается нормативное значение мощности класса.

В кроссовой комнате выделены две независимые подсистемы дистанционного питания.

\textbf{Подсистема питания видеокамер и СКУД (коммутатор Eltex MES2448P, PoE+)}

Оконечное оборудование распределено по следующим классам мощности:
\begin{itemize}
    \item Узконаправленные камеры QTECH QVC-MiR404Z532 (6 шт.): IEEE 802.3at (PoE+, Класс 4). Нормативная мощность порта: 30 Вт.
    \item Обзорные камеры QVC-MiR803M (2 шт.): IEEE 802.3at (PoE+, Класс 4). Нормативная мощность порта: 30 Вт.
    \item Контроллеры СКУД Parsec NC-8000-I (2 шт.): IEEE 802.3af (PoE, Класс 3). Нормативная мощность порта: 15,4 Вт.
\end{itemize}

Суммарный нормативный бюджет для данной группы портов составляет:
\begin{equation}
    P_{\text{PoE+}} = (6 + 2) \times 30 + 2 \times 15{,}4 = 240 + 30{,}8 = 270{,}8 \text{ Вт}.
\end{equation}

Аппаратный PoE-бюджет выбранного коммутатора Eltex MES2448P составляет 720 Вт. Данное значение с многократным избытком перекрывает расчетный классовый бюджет в 271 Вт, гарантируя стабильную подачу питания на камеры и устройства СКУД.

\textbf{Подсистема питания моноблоков операторов (коммутатор Aquarius, PoE++)}

Питание 6 моноблоков «Гравитон» М75И осуществляется через PoE-сплиттеры Planet POE-173S, подключенные к высокомощным портам коммутатора Aquarius AQ-N3000-48P4Y2Q. Сплиттеры классифицируются по максимальному уровню потребления:
\begin{itemize}
    \item PoE-сплиттер Planet POE-173S (6 шт.): IEEE 802.3bt Type 4 (PoE++, Класс 8). Нормативная мощность порта: 90 Вт.
\end{itemize}

Суммарный нормативный бюджет для рабочих мест операторов:
\begin{equation}
    P_{\text{PoE++}} = 6 \times 90 \text{ Вт} = 540 \text{ Вт}.
\end{equation}

Номинальный базовый PoE-бюджет коммутатора Aquarius составляет 1080 Вт. Устройство полностью покрывает расчетную нормативную нагрузку (540 Вт) с двукратным запасом, обеспечивая высочайшую надежность функционирования терминалов.

\chapter{Шкафы, ИБП, АКБ, охлаждение, компоновка}

\section{Электропитание и выбор источников бесперебойного питания}

Оборудование серверной и кроссовой комнат подключается к выделенным источникам бесперебойного питания (ИБП) с топологией On-Line (двойное преобразование) для обеспечения автономной работы не менее 30 минут. Расчет ведется по типичному энергопотреблению оборудования, а не по номиналам блоков питания.

\textbf{Расчет для телекоммуникационного шкафа серверной (пом. №\,19)}

В серверной располагается вычислительное оборудование без портов PoE. Типичное потребление:
\begin{itemize}
    \item Сервисный маршрутизатор Eltex ESR-1200: 60 Вт;
    \item Коммутатор ядра Eltex MES5310-48: 150 Вт;
    \item Сервер CI/CD Гравитон С2122АА: 300 Вт;
    \item Медиасервер Гравитон С2242И: 350 Вт;
    \item NAS Qtech QSRV-260202: 250 Вт.
\end{itemize}

Общее потребление оборудования (Потребление от сети):
\begin{equation}
    P = 60 + 150 + 300 + 350 + 250 = 1110 \text{ Вт}
\end{equation}

Ток нагрузки для выбора PDU (при напряжении $U = 230$ В и коэффициенте мощности $\cos \varphi = 0{,}95$):
\begin{equation}
    I = \frac{1110}{230 \cdot 0{,}95} \approx 5{,}1 \text{ А}
\end{equation}
Для серверной достаточно базового PDU на 16 А.

Требуемая активная мощность ИБП с запасом 30\,\%:
\begin{equation}
    P_{\text{ИБП}} \ge 1{,}3 \cdot 1110 \approx 1443 \text{ Вт}
\end{equation}
Выбирается стоечный ИБП \textbf{Сайбер Электро РСК-Эксперт-II-2000С} (2000 ВА / 1800 Вт). Загрузка составит $1110 / 1800 \approx 61{,}6\,\%$.

\textbf{Расчет для телекоммуникационного шкафа кроссовой (пом. №\,16)}

Для кроссовой учитывается собственное потребление коммутаторов и мощность PoE.
Состав:
\begin{itemize}
    \item PoE++ коммутатор Aquarius: собственное потребление 60 Вт;
    \item PoE+ коммутатор Eltex MES2448P: собственное потребление 40 Вт;
    \item Базовые коммутаторы Eltex MES2448B (2 шт.): $2 \times 40 = 80$ Вт.
\end{itemize}
Суммарная потребляемая мощность PoE-портов (по классам): $P_{\text{PoE}} = 540 + 270{,}8 = 810{,}8$ Вт.
Принимаем КПД блоков питания PoE-коммутаторов равным 0,9. Потребление от сети:
\begin{equation}
    P = \frac{60 + 40 + 810{,}8}{0{,}9} + 80 = 1012 + 80 = 1092 \text{ Вт}
\end{equation}

Ток нагрузки:
\begin{equation}
    I = \frac{1092}{230 \cdot 0{,}95} \approx 5{,}0 \text{ А}
\end{equation}
Достаточно PDU 16 А.

Требуемая активная мощность ИБП:
\begin{equation}
    P_{\text{ИБП}} \ge 1{,}3 \cdot 1092 \approx 1420 \text{ Вт}
\end{equation}
Выбирается ИБП \textbf{Сайбер Электро РСК-Эксперт-II-2000С} (2000 ВА / 1800 Вт). Загрузка $\approx 60\,\%$.

\section{Расчет емкости аккумуляторных батарей (АКБ)}

Время автономной работы системы должно составлять не менее 30 минут. При нагрузках 1110 Вт и 1092 Вт штатных внутренних батарей ИБП (4 шт. 12В/9Ач) хватает на $\approx 10$ минут (согласно таблицам автономии). Для обеспечения 30 минут автономности к каждому ИБП добавляется по одному внешнему стоечному батарейному модулю \textbf{БМ-РСК-Эксперт-II-2000} (2U, 8 шт. 12В/9Ач), что гарантирует работу в течение $\approx 35$ минут при указанной нагрузке.

\section{Тепловыделение в помещении}

Тепло в помещении — это потреблённая мощность минус мощность, ушедшая по кабелям за пределы помещения (PoE):
\begin{equation}
    Q = P_{\text{потр}} - P_{PoE,\ \text{отданная}} + \Delta P_{\text{ИБП}}
\end{equation}

\textbf{Серверная комната (пом. №\,19)}
1. Оборудование: 1110 Вт. PoE нет — всё потребление превращается в тепло.
2. Потери ИБП (с учетом нагрузки и зарядки батарей $P_{\text{вх}}$):
\begin{equation}
    1110 \cdot \left(\frac{1}{0{,}92} - 1\right) \approx 97 \text{ Вт}
\end{equation}
3. Освещение: $\approx 10$ Вт/м$^2 \times 10,7$ м$^2 = 107$ Вт. 
4. Итого тепловыделение серверной:
\begin{equation}
    Q_{\text{серв}} \approx 1110 + 97 + 107 = 1314 \text{ Вт} \approx 4484 \text{ BTU/ч}
\end{equation}

\textbf{Кроссовая комната (пом. №\,16)}
1. PoE-коммутаторы. Мощность PoE (810,8 Вт) уходит к конечным устройствам в другие помещения:
\begin{equation}
    Q_1 = 60 + 40 + 810{,}8 \cdot \left(\frac{1}{0{,}9} - 1\right) = 100 + 90 \approx 190 \text{ Вт}
\end{equation}
2. Коммутаторы без PoE:
\begin{equation}
    Q_2 = 80 \text{ Вт}
\end{equation}
3. Потери ИБП: Нагрузка ИБП составляет 1092 Вт.
\begin{equation}
    Q_3 = 1092 \cdot \left(\frac{1}{0{,}92} - 1\right) \approx 95 \text{ Вт}
\end{equation}
4. Освещение: $\approx 10$ Вт/м$^2 \times 11,7$ м$^2 = 117$ Вт.
5. Итого:
\begin{equation}
    Q_{\text{кросс}} = 190 + 80 + 95 + 117 \approx 482 \text{ Вт} \approx 1644 \text{ BTU/ч}
\end{equation}

Для ИБП в кроссовой нагрузка 1092 Вт, а тепла в шкафу около 365 Вт.

\section{Кондиционирование}

Требуемая холодопроизводительность с запасом 30\,\%:
\begin{equation}
    Q_{\text{хол}} = 1{,}3 \cdot Q
\end{equation}

Для серверной:
\begin{equation}
    Q_{\text{хол\_серв}} = 1{,}3 \cdot 1314 \approx 1708 \text{ Вт} \approx 1{,}7 \text{ кВт}
\end{equation}

Для кроссовой:
\begin{equation}
    Q_{\text{хол\_кросс}} = 1{,}3 \cdot 482 \approx 627 \text{ Вт} \approx 0{,}6 \text{ кВт}
\end{equation}

Для обоих помещений выбираются настенные сплит-системы с холодопроизводительностью 2,5 кВт (например, Dahatsu DHP-09). В каждом помещении устанавливаются по два кондиционера с блоком ротации для обеспечения 100\,\% резервирования и охлаждения по схеме 1+1. Кондиционеры питаются от отдельных линий электроснабжения, а не от ИБП шкафов.

\section{Вентиляция шкафа}

Требуемый расход воздуха через шкаф:
\begin{equation}
    L = \frac{3600 \cdot Q}{\rho \cdot c \cdot \Delta T}
\end{equation}
где $\rho = 1{,}2$ кг/м$^3$ — плотность воздуха, $c = 1005$ Дж/(кг$\cdot$К) — теплоемкость воздуха, $\Delta T = 10$ К — допустимый нагрев воздуха.

Серверный шкаф (оборудование и ИБП выделяют $1110 + 97 = 1207$ Вт тепла в шкафу):
\begin{equation}
    L_{\text{серв}} = \frac{3600 \cdot 1207}{1{,}2 \cdot 1005 \cdot 10} \approx 360 \text{ м}^3/\text{ч}
\end{equation}

Шкаф кроссовой (оборудование и ИБП выделяют $190 + 80 + 95 = 365$ Вт тепла в шкафу):
\begin{equation}
    L_{\text{кросс}} = \frac{3600 \cdot 365}{1{,}2 \cdot 1005 \cdot 10} \approx 109 \text{ м}^3/\text{ч}
\end{equation}

Для обеспечения такого расхода используются шкафы с перфорированными дверями и потолочные вентиляторные модули. Свободные юниты между оборудованием закрываются заглушками (фальш-панелями).

\section{Весогабаритный расчет и схемы компоновки оборудования}

\textbf{1. Шкаф серверной (пом. №\,19)}
\begin{itemize}
    \item Оборудование: Маршрутизатор (1U), ядро (1U), сервер CI/CD (2U), медиасервер (2U), NAS (2U), ИБП (2U), батарейный модуль (2U), консоль (1U), 3 органайзера и кросс-панель (4U). Итого: 17 U.
    \item Высота шкафа: $U_{\text{шкаф}} \ge \frac{17}{1 - 0{,}3} \approx 24{,}2$. Для серверной выбран шкаф \textbf{42U}. Запас $K_{\text{зап}} = \frac{17}{42} \cdot 100\,\% \approx 40\,\%$.
    \item Глубина и ширина: Серверы имеют глубину до 750 мм. Принимается шкаф глубиной 1200 мм и шириной 600 мм (серия C3.RK426010.2P2P (ГИСП 3938495)).
    \item Масса: Масса оборудования $\approx 150$ кг, шкафа $\approx 120$ кг. Итого $\approx 270$ кг. Нагрузка на перекрытие:
    \begin{equation}
        p = \frac{m \cdot g}{A} = \frac{270 \cdot 9{,}81}{0{,}6 \cdot 1{,}2} \approx 3{,}68 \text{ кПа}
    \end{equation}
    Это значение находится на грани типовой нормы офисных зданий (4 кПа), поэтому рекомендуется применение распределительных подиумов.
\end{itemize}

\textbf{2. Шкаф кроссовой (пом. №\,16)}
\begin{itemize}
    \item Оборудование: 4 коммутатора (4U), 7 патч-панелей (7U), 7 органайзеров (7U), оптика (1U), ИБП (2U), БМ (2U). Итого: 23 U.
    \item Высота шкафа: $U_{\text{шкаф}} \ge \frac{23}{0{,}7} \approx 32{,}8$. Для кроссовой также выбран шкаф \textbf{42U}.
    \item Глубина и ширина: Выбрана ширина 800 мм для удобной укладки большого числа пучков кабеля, глубина 1000 мм (C3.RI428010.2P2P (ГИСП 5776698)). 
    \item Масса: Оборудование весит менее 100 кг. Установка допустима без подиума.
\end{itemize}

\textbf{3. Размещение шкафов в помещении}
Для шкафа кроссовой длина по оси обслуживания: $L = L_{\text{перед}} + D_{\text{шкаф}} + L_{\text{зад}} = 1{,}0 + 1{,}0 + 0{,}8 = 2{,}8$ м.
Для шкафа серверной: $L = 1{,}0 + 1{,}2 + 0{,}8 = 3{,}0$ м.
Помещения лаборатории соответствуют этим требованиям.

\textbf{4. Схемы заполнения шкафов по юнитам (Фронтальный вид)}

\textbf{Таблица 1. Серверный шкаф 42U (Фронтальный вид)}

\begin{small}
\begin{longtable}{|c|l|l|}
\hline
\textbf{U} & \textbf{Оборудование} & \textbf{Примечание} \\ \hline
\endfirsthead
\hline
\textbf{U} & \textbf{Оборудование} & \textbf{Примечание} \\ \hline
\endhead
42 & Оптическая патч-панель & Ввод оптики от кроссовой \\ \hline
41 & Кабельный органайзер 1U & \\ \hline
40 & Маршрутизатор Eltex ESR-1200 & \\ \hline
39 & Кабельный органайзер 1U & \\ \hline
38 & Коммутатор ядра Eltex MES5324 & \\ \hline
37 & Кабельный органайзер 1U & \\ \hline
36--22 & Резерв (Фальш-панели) & 15U \\ \hline
21 & Консоль KVM 1U & На удобной высоте $\sim$1 м \\ \hline
20--19 & Резерв (Фальш-панели) & 2U \\ \hline
18--17 & Сервер CI/CD Гравитон С2122АА & Вычислительное оборудование \\ \hline
16--15 & Медиасервер Гравитон С2242И & Вычислительное оборудование \\ \hline
14--13 & NAS Qtech QSRV-260202 & Массивные диски \\ \hline
12--5 & Резерв (Фальш-панели) & 8U под будущее расширение \\ \hline
4--3 & ИБП 2000 ВА & Источник питания \\ \hline
2--1 & Батарейный модуль 2U & В самом низу (Тяжелое) \\ \hline
0 & Вертикальный PDU 16 А & На задней стойке \\ \hline
\end{longtable}
\end{small}

\textbf{Таблица 2. Шкаф кроссовой 42U (Фронтальный вид)}

\begin{small}
\begin{longtable}{|c|l|l|}
\hline
\textbf{U} & \textbf{Оборудование} & \textbf{Примечание} \\ \hline
\endfirsthead
\hline
\textbf{U} & \textbf{Оборудование} & \textbf{Примечание} \\ \hline
\endhead
42 & Оптическая патч-панель & Аплинк до серверной \\ \hline
41 & Кабельный органайзер 1U & \\ \hline
40 & Патч-панель 24xRJ45 & Для коммутатора Aquarius \\ \hline
39 & Коммутатор Aquarius PoE++ & \\ \hline
38 & Патч-панель 24xRJ45 & Для коммутатора Aquarius \\ \hline
37 & Кабельный органайзер 1U & \\ \hline
36 & Патч-панель 24xRJ45 & Для Eltex PoE+ \\ \hline
35 & Коммутатор Eltex MES2448P & \\ \hline
34 & Патч-панель 24xRJ45 & Для Eltex PoE+ \\ \hline
33 & Кабельный органайзер 1U & \\ \hline
32 & Патч-панель 24xRJ45 & Для Eltex MES2448B (1) \\ \hline
31 & Коммутатор Eltex MES2448B (1) & \\ \hline
30 & Патч-панель 24xRJ45 & Для Eltex MES2448B (1) \\ \hline
29 & Кабельный органайзер 1U & \\ \hline
28 & Патч-панель 24xRJ45 & Для Eltex MES2448B (2) \\ \hline
27 & Коммутатор Eltex MES2448B (2) & \\ \hline
26 & Патч-панель 24xRJ45 & Для Eltex MES2448B (2) \\ \hline
25 & Кабельный органайзер 1U & \\ \hline
24--5 & Резерв (Фальш-панели) & 20U \\ \hline
4--3 & ИБП 2000 ВА & \\ \hline
2--1 & Батарейный модуль 2U & \\ \hline
0 & Вертикальный PDU 16 А & На задней стойке \\ \hline
\end{longtable}
\end{small}

\chapter{Расчет длины кабелей горизонтальной подсистемы}
(В разработке)

\chapter{Емкость кабельных трасс и заполнение лотков}
(В разработке)

\chapter{Кабельный журнал}
(В разработке)

\chapter{Экономическая часть}
(В разработке)

\chapter{Безопасность жизнедеятельности}
(В разработке)

\cleardoublepage
\phantomsection
\addcontentsline{toc}{chapter}{ЗАКЛЮЧЕНИЕ}
\chapter*{ЗАКЛЮЧЕНИЕ}

В ходе выполнения курсового проекта была спроектирована структурированная кабельная система (СКС) этажа лаборатории тестирования устройств IoT. Архитектура сети построена по топологии иерархической звезды в соответствии с современными стандартами отказоустойчивости.

В рамках работы были решены следующие основные задачи:
\begin{enumerate}
    \item Произведен анализ планировки этажа и рассчитано необходимое количество информационных портов, составившее 148 портов. Учтены специфические требования лабораторных стендов с использованием многопользовательских розеток (MUTOA).
    \item Выполнен расчет видеотрафика от 8 IP-камер с учетом применения кодеков H.265. Подтверждена достаточность пропускной способности каналов Gigabit Ethernet на уровне доступа и необходимость агрегированных оптических каналов 10G SFP+ на магистральных участках.
    \item Обоснован выбор активного сетевого оборудования из Единого реестра российской радиоэлектронной продукции Минпромторга РФ. Подобраны коммутаторы Eltex и Aquarius, обеспечивающие необходимый бюджет мощности (PoE+ и PoE++) для питания камер, СКУД и моноблоков операторов.
    \item Выполнен расчет емкости и заполняемости телекоммуникационных шкафов. Спроектированы кабельные трассы и произведена оценка тепловыделения активного оборудования, на основании которой подобраны системы кондиционирования и источники бесперебойного питания (ИБП).
    \item Разработан комплект технической документации, включающий техническое задание по ГОСТ 34.602-2020, кабельный журнал, ведомость покупных изделий по форме ГОСТ 2.106 и чертеж схемы расположения оборудования.
\end{enumerate}

Спроектированная СКС обладает высокой надежностью, соответствует требованиям к пропускной способности при пиковых нагрузках, обеспечивает питание конечных устройств по технологии PoE и имеет необходимый резерв для будущего масштабирования лаборатории. Требования технического задания выполнены в полном объеме.
\cleardoublepage
\phantomsection
\addcontentsline{toc}{chapter}{СПИСОК ИСПОЛЬЗОВАННЫХ ИСТОЧНИКОВ}
\chapter*{СПИСОК ИСПОЛЬЗОВАННЫХ ИСТОЧНИКОВ}

\begin{enumerate}
    \item ГОСТ 34.602-2020. Информационные технологии (ИТ). Комплекс стандартов на автоматизированные системы. Техническое задание на создание автоматизированной системы. -- Введ. 2022-01-01. -- Москва : Стандартинформ, 2021. -- 18 с.
    \item ГОСТ Р 15.301-2016. Система разработки и постановки продукции на производство. Продукция производственно-технического назначения. Порядок разработки и постановки продукции на производство. -- Введ. 2017-07-01. -- Москва : Стандартинформ, 2017. -- 18 с.
    \item ГОСТ Р 53246-2025. Информационные технологии. Системы кабельные структурированные. Проектирование основных узлов системы. Общие требования. -- Введ. 2026-01-01. -- Москва : Российский институт стандартизации, 2025. -- 82 с.
    \item ГОСТ Р 53245-2008. Информационные технологии. Системы кабельные структурированные. Монтаж основных узлов системы. Методы испытания. -- Введ. 2010-01-01. -- Москва : Стандартинформ, 2009. -- 20 с.
    \item ГОСТ Р 50571.5.54-2013/МЭК 60364-5-54:2011. Электроустановки низковольтные. Часть 5-54. Выбор и монтаж электрооборудования. Заземляющие устройства, защитные проводники и защитные проводники уравнивания потенциалов. -- Введ. 2015-01-01. -- Москва : Стандартинформ, 2014. -- 36 с.
    \item ГОСТ 31565-2012. Кабельные изделия. Требования пожарной безопасности. -- Введ. 2014-01-01. -- Москва : Стандартинформ, 2013. -- 10 с.
    \item ГОСТ Р 70443-2022. Информационные технологии (ИТ). Системы кабельные структурированные. Техническое задание. Требования к составу и содержанию. -- Введ. 2023-01-01. -- Москва : Российский институт стандартизации, 2022. -- 16 с.
    \item ГОСТ 2.106-2019. Единая система конструкторской документации. Текстовые документы. -- Введ. 2020-02-01. -- Москва : Стандартинформ, 2019. -- 35 с.
    \item Государственная информационная система промышленности (ГИСП). Каталог промышленной продукции. -- URL: \url{https://gisp.gov.ru/goods/} (дата обращения: 03.10.2026).
\end{enumerate}

\cleardoublepage
\begin{appendices}
\stepcounter{chapter}
\phantomsection
\addcontentsline{toc}{chapter}{ПРИЛОЖЕНИЕ А}

% =======================================================================
%  ТИТУЛЬНАЯ СТРАНИЦА ТЗ
% =======================================================================
\begin{titlepage}
    \thispagestyle{empty}
    \begin{flushright}
    ПРИЛОЖЕНИЕ А
    \end{flushright}

    % --- Верхний блок: СОГЛАСОВАНО / УТВЕРЖДАЮ ---
    \vspace*{0pt}
    \noindent
    \begin{minipage}[t]{0.48\textwidth}
        \raggedright
        \textbf{СОГЛАСОВАНО}\\[6pt]
        Руководитель НИР\\[10pt]
        \blank[3cm]~***\\[6pt]
        <<\blank[0.6cm]>>~\blank[2.5cm]~2026~г.
    \end{minipage}%
    \hfill
    \begin{minipage}[t]{0.48\textwidth}
        \raggedright
        \textbf{УТВЕРЖДАЮ}\\[6pt]
        Заведующий кафедрой ИУ3\\[10pt]
        \blank[3cm]~***\\[6pt]
        <<\blank[0.6cm]>>~\blank[2.5cm]~2026~г.
    \end{minipage}

    % --- Центральный блок: Название ТЗ и тема ---
    \vfill
    \begin{center}
        {\LARGE \bfseries ТЕХНИЧЕСКОЕ ЗАДАНИЕ НА НИР}\\[8mm]
        {\Large \bfseries
        ПРОЕКТИРОВАНИЕ СТРУКТУРИРОВАННОЙ КАБЕЛЬНОЙ\\
        СИСТЕМЫ ОФИСА ЛАБОРАТОРИИ\\
        ТЕСТИРОВАНИЯ УСТРОЙСТВ IoT}
    \end{center}
    \vfill

    % --- Нижний блок: РАЗРАБОТАНО ---
    \noindent
    \textbf{РАЗРАБОТАНО}\\[6pt]
    Студент *\\[10pt]
    \blank[3cm]~***\\[6pt]
    <<\blank[0.6cm]>>~\blank[2.5cm]~2026~г.

    \vspace{20mm}

    % --- Город и год ---
    \begin{center}
        Москва\\
        2026~г.
    \end{center}
\end{titlepage}

\setcounter{page}{2}

% =======================================================================
%  РАЗДЕЛ 1. ОБЩИЕ СВЕДЕНИЯ
% =======================================================================

\addtocontents{toc}{\protect\setcounter{tocdepth}{0}}
\renewcommand{\thesection}{\arabic{section}}
\setcounter{section}{0}
\section{ОБЩИЕ СВЕДЕНИЯ}

\subsection{Полное наименование системы и ее условное обозначение}
\begin{enumerate}[label=\thesection.\arabic*]
    \item \textbf{Полное наименование:} Структурированная кабельная система (СКС) этажного офиса лаборатории тестирования аппаратуры и программного обеспечения устройств IoT.
    \item \textbf{Условное обозначение:} СКС-IoT-ИУ3-V10.
\end{enumerate}

\subsection{Шифр темы или основание для проведения работ}
\begin{enumerate}[label=\thesection.\arabic*]
    \item \textbf{Основание для разработки:} Индивидуальное задание к курсовой работе по дисциплине <<Проектирование структурированной кабельной системы>>, 3-й семестр бакалавриата кафедры ИУ-3.
\end{enumerate}

\subsection{Наименование предприятий (организаций) --- Заказчика и Разработчика}
\begin{enumerate}[label=\thesection.\arabic*]
    \item \textbf{Заказчик:} Кафедра ИУ-3 МГТУ им. Н.Э. Баумана.
    \item \textbf{Разработчик:} Студент группы *, ***.
\end{enumerate}

\subsection{Сроки начала и окончания работы}
\begin{enumerate}[label=\thesection.\arabic*]
    \item \textbf{Начало:} Сентябрь 2026 г.
    \item \textbf{Окончание:} Декабрь 2026 г.
\end{enumerate}

% =======================================================================
%  РАЗДЕЛ 2. ЦЕЛИ И НАЗНАЧЕНИЕ СОЗДАНИЯ АВТОМАТИЗИРОВАННОЙ СИСТЕМЫ
% =======================================================================

\section{ЦЕЛИ И НАЗНАЧЕНИЕ СОЗДАНИЯ КАБЕЛЬНОЙ СИСТЕМЫ}

\subsection{Назначение системы}
СКС предназначена для создания универсальной, надежной и масштабируемой кабельной инфраструктуры, обеспечивающей передачу данных, сигналов управления и электропитания по технологии PoE (Power over Ethernet) для обеспечения функционирования:
\begin{enumerate}[label=\asbuk*)]
    \item лабораторных стендов тестирования IoT-устройств и контрольно-измерительного оборудования;
    \item рабочих мест операторов и административного персонала;
    \item систем охранного видеонаблюдения контроля стендов и общих зон;
    \item систем контроля и управления доступом (СКУД);
    \item магистральных каналов связи с серверными здания.
\end{enumerate}

\subsection{Цели создания системы}
\begin{enumerate}[label=\thesection.\arabic*]
    \item \textbf{Обеспечение требуемой пропускной способности} для передачи трафика от IoT-устройств и систем видеонаблюдения высокого разрешения (в режимах MCU и SFU).
    \item \textbf{Повышение надежности} за счет организации топологии <<иерархическая звезда>> и избыточности магистральных линий.
    \item \textbf{Оптимизация эксплуатации} ИТ-инфраструктуры за счет сведения всех кабельных систем к единому коммутационному центру.
    \item \textbf{Выполнение импортозамещения} путем подбора активного сетевого оборудования из Единого реестра российской радиоэлектронной продукции Минпромторга РФ.
\end{enumerate}

% =======================================================================
%  РАЗДЕЛ 3. ХАРАКТЕРИСТИКА КАБЕЛЬНОЙ СИСТЕМЫ
% =======================================================================

\section{ХАРАКТЕРИСТИКА КАБЕЛЬНОЙ СИСТЕМЫ}

\subsection{Краткие сведения об объекте}
\begin{enumerate}[label=\thesection.\arabic*]
    \item \textbf{Расположение:} Один этаж в многоэтажном здании.
    \item \textbf{Общая площадь:} 613,08~м$^2$ в осях.
    \item \textbf{Количество помещений:} 24 (согласно экспликации архитектурного плана формата А2).
    \item \textbf{Режим работы:} Непрерывный, с повышенной нагрузкой на кабельные тракты в зонах испытаний.
\end{enumerate}

\subsection{Перечень оборудуемых помещений и плотность рабочих мест}
Распределение информационных портов и оборудования выполняется строго по экспликации:

\begin{enumerate}[label=\thesection.\arabic*]
    \item \textbf{Помещения лабораторий:}
    \begin{enumerate}[label=\thesection.\arabic{enumi}.\arabic*]
        \item \textbf{Помещение №4 (Кабинет 3):} 3~лабораторных стенда. На каждом стенде: 18~портов IoT-устройств, 1~порт камеры контроля (PoE+), 6~резервных портов. Емкость MUTOA стенда = 25~портов.
        \item \textbf{Помещение №11 (Кабинет 1):} 3~лабораторных стенда. На каждом стенде: 8~портов IoT-устройств, 1~порт камеры контроля (PoE+), 6~резервных портов. Емкость MUTOA стенда = 15~портов.
        \item \textit{Дополнительно в №4 и №11:} по 1~камере общего вида (PoE+) и по 1~контроллеру СКУД (PoE+).
    \end{enumerate}

    \item \textbf{Операторская (Помещение №20 <<Опенспейс>>):}
    \begin{enumerate}[label=\thesection.\arabic{enumi}.\arabic*]
        \item \textbf{Пост №1} (обслуживает 1~стенд): 1~моноблок (1~порт PoE++), 1~рабочая станция (1~порт), 1~резервный порт. Итого: 3~розетки.
        \item \textbf{Пост №2} (обслуживает 1~стенд): 1~моноблок (1~порт PoE++), 1~рабочая станция (1~порт), 1~резервный порт. Итого: 3~розетки.
        \item \textbf{Пост №3} (обслуживает 1~стенд): 1~моноблок (1~порт PoE++), 1~рабочая станция (1~порт), 1~резервный порт. Итого: 3~розетки.
        \item \textbf{Пост №4} (обслуживает 1~стенд): 1~моноблок (1~порт PoE++), 1~рабочая станция (1~порт), 1~резервный порт. Итого: 3~розетки.
        \item \textbf{Пост №5} (обслуживает 1~стенд): 1~моноблок (1~порт PoE++), 1~рабочая станция (1~порт), 1~резервный порт. Итого: 3~розетки.
        \item \textbf{Пост №6} (обслуживает 1~стенд): 1~моноблок (1~порт PoE++), 1~рабочая станция (1~порт), 1~резервный порт. Итого: 3~розетки.
        \item \textbf{Общая зона опенспейса:} 2~розетки СКС и 1~штатное место персонала (по 2~розетки).
    \end{enumerate}

    \item \textbf{Административные помещения:}
    \begin{enumerate}[label=\thesection.\arabic{enumi}.\arabic*]
        \item Помещения №24 (Кабинет 2), №7 (Переговорная): оборудуются 2~розетками СКС каждое.
    \end{enumerate}

    \item \textbf{Технологические помещения:}
    \begin{enumerate}[label=\thesection.\arabic{enumi}.\arabic*]
        \item \textbf{Помещение №19:} Серверная этажа (11,70~м$^2$). Размещение главного шкафа.
        \item \textbf{Помещение №16:} Кроссовая СКС (11,70~м$^2$). Размещение этажного распределительного пункта.
    \end{enumerate}
\end{enumerate}

% =======================================================================
%  РАЗДЕЛ 4. ТРЕБОВАНИЯ К АВТОМАТИЗИРОВАННОЙ СИСТЕМЕ
% =======================================================================

\section{ТРЕБОВАНИЯ К КАБЕЛЬНОЙ СИСТЕМЕ}

\subsection{Требования к структуре и архитектуре системы}
\begin{enumerate}[label=\thesection.\arabic*]
    \item СКС должна быть спроектирована по архитектуре \textbf{иерархической звезды} с одним техническим центром (Серверная №19) и одной кроссовой комнатой (Помещение №16).
    \item В кроссовой СКС (№16) должен быть организован Этажный Распределительный Пункт (ЭРП).
    \item Магистральная подсистема должна связывать ЭРП этажа со слаботочным стояком (<<СКС>>) и организовывать вертикальные связи с Главной серверной здания (пом. 903, эт. 9) и Резервной серверной (пом. 1183, эт. 11).
\end{enumerate}

\subsection{Требования к производительности и компонентам}
\begin{enumerate}[label=\thesection.\arabic*]
    \item Кабельная система горизонтальной подсистемы должна быть выполнена на основе медного четырехпарного кабеля <<витая пара>> категории не ниже \textbf{Cat.6A (класс Ea)} для поддержки скоростей до 10~Гбит/с и обеспечения стабильной работы PoE++/4PPoE (IEEE 802.3bt).
    \item Магистральная подсистема здания (вертикальная) должна быть спроектирована с использованием волоконно-оптического многомодового (OM4) или одномодового (OS2) кабеля для обеспечения отказоустойчивости и емкости каналов.
    \item Коммутационные панели (патч-панели) в шкафах должны иметь плотность не менее 24 портов RJ-45/1U с обязательной цветовой или цифровой маркировкой подсистем.
\end{enumerate}

\subsection{Требования к надежности и режимам видеонаблюдения}
Система должна обеспечивать бесперебойную передачу трафика от 8~камер (6~направленных, 2~обзорных) с поддержкой кодеков H.265/H.265+ в двух режимах работы медиасерверов:
\begin{enumerate}[label=\asbuk*)]
    \item штатный режим (MCU) --- архитектура с централизованным микшированием потоков;
    \item аварийный режим (SFU) --- архитектура с маршрутизацией сырых потоков напрямую на операторов, требующая проверки пиковых пропускных способностей каналов операторской подсистемы.
\end{enumerate}

\subsection{Требования к условиям эксплуатации и климатике}
\begin{enumerate}[label=\thesection.\arabic*]
    \item В помещениях Серверной (№19) и Кроссовой (№16) должна быть предусмотрена принудительно-вытяжная вентиляция или кондиционирование для компенсации суммарного тепловыделения активного оборудования.
    \item Параметры микроклимата: температура 18--24~$^\circ$C, относительная влажность 30--55\% без конденсации.
    \item Active-оборудование шкафов должно питаться от выделенных источников бесперебойного питания (ИБП) с топологией On-Line и временем автономной работы при максимальной нагрузке не менее 30 минут.
\end{enumerate}

\subsection{Требования к кабельным трассам и конструкциям}
\begin{enumerate}[label=\thesection.\arabic*]
    \item Магистральные лотки в коридорах должны быть металлическими (проволочными или перфорированными).
    \item Заполнение лотков и кабельных каналов не должно превышать \textbf{40--50}\% на этапе проектирования с учетом требований по радиусу изгиба и перспективного расширения.
    \item Расчет подвесов и элементов крепления должен проводиться с коэффициентом запаса прочности не менее 3 по отношению к весу заполненного лотка.
\end{enumerate}

\subsection{Требования к заземлению и электробезопасности}
\begin{enumerate}[label=\thesection.\arabic*]
    \item Все металлические лотки, корпуса шкафов (19"), патч-панели и экраны кабелей должны быть интегрированы в единый контур защитного технологического заземления здания согласно ГОСТ Р 50571.5.54.
    \item Кабели должны отвечать требованиям пожарной безопасности по \textbf{ГОСТ 31565-2012} (исполнение не ниже нг-HF или нг-FRHF).
\end{enumerate}

% =======================================================================
%  РАЗДЕЛ 5. СОСТАВ И СОДЕРЖАНИЕ РАБОТ ПО СОЗДАНИЮ КС
% =======================================================================

\section{СОСТАВ И СОДЕРЖАНИЕ РАБОТ ПО СОЗДАНИЮ КАБЕЛЬНОЙ СИСТЕМЫ}
В рамках курсового проектирования (согласно стадиям разработки по ГОСТ Р 15.301-2016) выполняются следующие этапы.
\begin{enumerate}[label=\thechapter.\arabic*]
    \item Анализ исходных данных. Изучение DXF-плана, расчет точного числа портов СКС для каждого помещения.
    \item Расчетно-теоретический этап:
    \begin{enumerate}[label=\asbuk*)]
        \item трассировка и расчет длин кабельных линий (метод контрольных длин);
        \item расчет сечения коробов, лотков, нагрузок на несущие конструкции;
        \item проектирование активных компонентов по реестру Минпромторга (коммутаторы доступа с PoE+/PoE++, агрегаторы, маршрутизаторы);
        \item расчет тепловыделения шкафов, параметров охлаждения и емкости аккумуляторных батарей ИБП.
    \end{enumerate}
    \item Графический этап. Создание схемы размещения рабочих мест на архитектурном плане (А2), составление схемы логической топологии и компоновки шкафов.
\end{enumerate}


% =======================================================================
%  РАЗДЕЛ 6. ПОРЯДОК КОНТРОЛЯ И ПРИЕМКИ КАБЕЛЬНОЙ СИСТЕМЫ
% =======================================================================

\section{ПОРЯДОК КОНТРОЛЯ И ПРИЕМКИ КАБЕЛЬНОЙ СИСТЕМЫ}

Контроль результатов учебной разработки осуществляет руководитель НИР.

При контроле проверяются:

\begin{itemize}[label=\asbuk*)]
    \item соответствие проектного решения исходным данным индивидуального варианта;
    \item корректность размещения оборудования и кабельных трасс;
    \item корректность расчета количества портов, патч-панелей и заполнения шкафов;
    \item корректность расчетов трафика в режимах MCU и SFU;
    \item корректность расчетов кабельных линий, лотков, подвесов, охлаждения и ИБП;
    \item обоснованность выбора активного оборудования и PoE-бюджета;
    \item соответствие текстовой и графической документации требованиям задания и применимых нормативных документов.
\end{itemize}

В рамках учебной НИР фактический монтаж и натурные приемочные испытания не выполняются. При реализации проекта кабельные линии подлежат испытаниям и сертификационным измерениям в соответствии с требованиями выбранной кабельной системы и применимых стандартов.

% =======================================================================
%  РАЗДЕЛ 7. ТРЕБОВАНИЯ К СОСТАВУ И СОДЕРЖАНИЮ РАБОТ ПО ПОДГОТОВКЕ
%            ОБЪЕКТА АВТОМАТИЗАЦИИ К ВВОДУ КС В ДЕЙСТВИЕ
% =======================================================================

\section{ТРЕБОВАНИЯ К СОСТАВУ И СОДЕРЖАНИЮ РАБОТ ПО ПОДГОТОВКЕ ОБЪЕКТА АВТОМАТИЗАЦИИ К ВВОДУ КАБЕЛЬНОЙ СИСТЕМЫ В ДЕЙСТВИЕ}

Перед вводом кабельной системы в эксплуатацию должны быть выполнены:

\begin{enumerate}[label=\asbuk*)]
    \item подготовка серверной №19 и кроссовой №16 к размещению телекоммуникационных шкафов.
    \item подготовка мест прокладки лотков, коробов и кабельных трасс.
    \item обеспечение электроснабжения, защитного заземления и охлаждения.
    \item монтаж и маркировка информационных розеток, MUTOA, патч-панелей и кабельных линий.
    \item установка и настройка активного оборудования.
    \item проверка кабельных линий и сетевых соединений.
    \item проверка работы PoE+/PoE++ оборудования и резервного питания.
    \item оформление исполнительной документации.
\end{enumerate}

% =======================================================================
%  РАЗДЕЛ 8. ТРЕБОВАНИЯ К ДОКУМЕНТИРОВАНИЮ
% =======================================================================

\section{ТРЕБОВАНИЯ К ДОКУМЕНТИРОВАНИЮ}

По результатам НИР должны быть подготовлены:

\begin{enumerate}[label=\asbuk*)]
    \item расчётно-пояснительная записка объёмом от 60 до 100 листов;
    \item техническое задание на кабельную систему;
    \item чертёж схемы сети объекта на листе формата A2;
    \item схема логической топологии сети;
    \item ведомость покупных изделий;
    \item расчёты портов, патч-панелей, шкафов, трафика, кабельных трасс, охлаждения и ИБП.
\end{enumerate}

% =======================================================================
%  РАЗДЕЛ 9. ИСТОЧНИКИ РАЗРАБОТКИ
% =======================================================================

\section{ИСТОЧНИКИ РАЗРАБОТКИ}

При разработке используются:

\begin{enumerate}[label=\asbuk*)]
    \item индивидуальное задание на выполнение НИР по теме <<Проектирование структурированной кабельной системы офиса лаборатории тестирования устройств IoT>>;
    \item исходные данные и архитектурный план индивидуального варианта;
    \item DXF-подоснова индивидуального варианта;
    \item ГОСТ Р 70443-2022, устанавливающий требования к техническому заданию на кабельную систему;
    \item ГОСТ 34.602-2020 и ГОСТ Р 15.301-2016, указанные в индивидуальном задании;
    \item ГОСТ 7.32-2017 в части оформления расчетно-пояснительной записки;
    \item ГОСТ Р 53246-2025 и ГОСТ Р 53245-2008 в части проектирования, монтажа и испытаний структурированных кабельных систем;
    \item ГОСТ Р 50571.5.54 в части защитного заземления;
    \item ГОСТ 31565-2012 в части пожарной безопасности кабельных изделий;
    \item Единый реестр российской радиоэлектронной продукции Минпромторга Российской Федерации.
\end{enumerate}




При разработке используются:

\begin{enumerate}[label=\asbuk*)]
    \item индивидуальное задание на выполнение НИР по теме <<Проектирование структурированной кабельной системы офиса лаборатории тестирования устройств IoT>>;
    \item исходные данные и архитектурный план индивидуального варианта;
    \item DXF-подоснова индивидуального варианта;
    \item ГОСТ Р 70443-2022, устанавливающий требования к техническому заданию на кабельную систему;
    \item ГОСТ 34.602-2020 и ГОСТ Р 15.301-2016, указанные в индивидуальном задании;
    \item ГОСТ 7.32-2017 в части оформления расчетно-пояснительной записки;
    \item ГОСТ Р 53246-2025 и ГОСТ Р 53245-2008 в части проектирования, монтажа и испытаний структурированных кабельных систем;
    \item ГОСТ Р 50571.5.54 в части защитного заземления;
    \item ГОСТ 31565-2012 в части пожарной безопасности кабельных изделий;
    \item Единый реестр российской радиоэлектронной продукции Минпромторга Российской Федерации.
\end{enumerate}

\vspace{1em}
\cleardoublepage
\stepcounter{chapter}
\phantomsection
\addtocontents{toc}{\protect\setcounter{tocdepth}{2}}
\renewcommand{\thesection}{\thechapter.\arabic{section}}
\addcontentsline{toc}{chapter}{ПРИЛОЖЕНИЕ Б}

\begin{flushright}
ПРИЛОЖЕНИЕ Б
\end{flushright}

\begin{center}
\textbf{\large ВЕДОМОСТЬ ПОКУПНЫХ ИЗДЕЛИЙ}
\end{center}

\begin{landscape}
\begin{small}
\begin{longtable}{|p{0.8cm}|p{2.5cm}|p{6.2cm}|p{2.3cm}|p{2.3cm}|p{3.2cm}|p{1.2cm}|p{3.0cm}|}
\hline
\centering \textbf{Поз.} &
\centering \textbf{Код продукции} &
\centering \textbf{Наименование} &
\centering \textbf{Обозначение документа} &
\centering \textbf{Поставщик} &
\centering \textbf{Куда входит} &
\centering \textbf{Кол-во} &
\centering \textbf{Примечание} \tabularnewline
\hline
\endfirsthead

\hline
\centering \textbf{Поз.} &
\centering \textbf{Код продукции} &
\centering \textbf{Наименование} &
\centering \textbf{Обозначение документа} &
\centering \textbf{Поставщик} &
\centering \textbf{Куда входит} &
\centering \textbf{Кол-во} &
\centering \textbf{Примечание} \tabularnewline
\hline
\endhead
\hline \multicolumn{8}{r}{\textit{Продолжение на следующей странице}} \\
\endfoot
\hline
\endlastfoot

\multicolumn{8}{|c|}{\textbf{Активное сетевое оборудование}} \\ \hline
1 & 4812300 & Коммутатор доступа AQ-N3000-48P4Y2Q & & Aquarius & Кроссовая & 1 & \\ \hline
2 & 3501245 & Коммутатор доступа MES2448B & & Eltex & Кроссовая & 2 & \\ \hline
3 & 3501246 & Коммутатор доступа MES2448P & & Eltex & Кроссовая & 1 & \\ \hline
4 & 5513820 & Коммутатор ядра MES5310-48 & & Eltex & Серверная & 1 & \\ \hline
6 & 5517235 & Сервисный маршрутизатор ESR-1200 & & Eltex & Серверная & 1 & \\ \hline
7 & Трансивер & Оптический трансивер SFP+ 10GBASE-SR & & Eltex & Коммутаторы & 4 & \\ \hline

\multicolumn{8}{|c|}{\textbf{Вычислительное оборудование и рабочие станции}} \\ \hline
8 & 4376653 & Медиасервер стоечный С2242И (2U) & & Гравитон & Серверная & 1 & \\ \hline
9 & 4533887 & Моноблок оператора М75И (PoE++) & & Гравитон & Операторская & 6 & \\ \hline
10 & 5404596 & Сервер CI/CD С2122АА (2U) & & Гравитон & Серверная & 1 & \\ \hline
11 & 2075661 & Сетевое хранилище (NAS) QSRV-260202 & & Qtech & Серверная & 1 & \\ \hline

\multicolumn{8}{|c|}{\textbf{Системы безопасности и контроля}} \\ \hline
12 & 4001025 & Обзорная 4 Мп IP-камера QVC-MiR803M & & QTECH & Лаборатории & 2 & \\ \hline
13 & 5531204 & Контроллер СКУД NC-8000-I (PoE) & & Parsec & Двери & 2 & \\ \hline
14 & 4001037 & Узконаправленная 4 Мп IP-камера QVC-MiR404Z532 & & QTECH & Стенды & 6 & \\ \hline

\multicolumn{8}{|c|}{\textbf{Системы бесперебойного электропитания и климат-контроля}} \\ \hline
15 & 5689102 & Батарейный модуль БМ-РСК-Эксперт-II-2000 & & Сайбер Электро & Шкафы & 2 & \\ \hline
16 & 5689101 & ИБП стоечный РСК-Эксперт-II-2000С & & Сайбер Электро & Шкафы & 2 & \\ \hline
17 & 3841120 & Блок распределения питания (PDU) 16А, 19'' & & C3 Solutions & Шкафы & 2 & \\ \hline
18 & Вне реестра & PoE-сплиттер Gigabit IEEE 802.3bt POE-173S & & Planet & Операторская & 6 & \\ \hline
19 & Вне реестра & Сплит-система настенная DHP-09 (2.5 кВт) & & Dahatsu & Серверная & 4 & \\ \hline
20 & 5776698 & Шкаф телекоммуникационный C3.RI (42U) & & C3 Solutions & Кроссовая & 1 & \\ \hline
21 & 3938495 & Шкаф телекоммуникационный C3.RK (42U) & & C3 Solutions & Серверная & 1 & \\ \hline
22 & Вентилятор & Модуль вентиляторный в крышу & & C3 Solutions & Шкафы & 2 & \\ \hline

\multicolumn{8}{|c|}{\textbf{Пассивное сетевое оборудование и кабельные системы}} \\ \hline
23 & NKL 4140C-BK & Кабель U/UTP, Cat.6A, 4x2x23AWG (бухта 500 м) & & NIKOMAX & Трассы & 8 & 4000 м \\ \hline
24 & NMC-RP24UD2 & Патч-панель 24 порта RJ-45, Cat.6A, 1U & & NIKOMAX & Шкафы & 7 & \\ \hline
25 & NMC-WP02UD2 & Розетка телекоммуникационная 2 порта RJ-45, Cat.6A & & NIKOMAX & Рабочие места & 75 & \\ \hline
26 & C3.ORG & Органайзер кабельный горизонтальный 19", 1U & & C3 Solutions & Шкафы & 10 & \\ \hline
27 & ОМ4-8 & Кабель волоконно-оптический многомодовый OM4 & & Еврокабель & Магистрали & 200 & В метрах \\ \hline
28 & Кросс 1U & Оптическая распределительная панель 19", 1U, 24 LC & & ЦМО & Шкафы & 2 & \\ \hline
29 & LC-LC & Шнур оптический соединительный LC-LC OM4, 2 м & & NIKOMAX & Шкафы & 4 & \\ \hline

\end{longtable}
\end{small}
\end{landscape}

\cleardoublepage
\stepcounter{chapter}
\phantomsection
\addcontentsline{toc}{chapter}{ПРИЛОЖЕНИЕ В}

\begin{flushright}
ПРИЛОЖЕНИЕ В
\end{flushright}

\begin{center}
\textbf{\large Схема логической топологии сети}
\end{center}

Ниже представлена структурная логическая схема разрабатываемой сети с разделением на уровни иерархии (Ядро, Доступ) и физическим расположением оборудования по помещениям.

\vspace{1cm}
\begin{figure}[H]
\centering
\resizebox{\textwidth}{!}{
\begin{tikzpicture}[
    box/.style={draw, rectangle, rounded corners, minimum width=3.5cm, minimum height=1cm, text centered, font=\sffamily\footnotesize},
    core/.style={box, fill=red!10, draw=red!60, thick},
    acc/.style={box, fill=green!10, draw=green!60, thick},
    server/.style={box, fill=gray!10, draw=gray!60},
    cloud/.style={draw, rectangle, rounded corners=0.5cm, fill=blue!5, draw=blue!40, thick, minimum width=3.5cm, minimum height=1cm, text centered, font=\sffamily\footnotesize},
    link/.style={draw, thick, -},
    uplink/.style={draw, thick, blue, -},
    node distance=2cm and 1cm
]

% Интернет
\node[cloud] (internet) {Интернет / Провайдер};

% Маршрутизатор (Шлюз)
\node[server, below=of internet, yshift=-0.5cm] (router) {Шлюз: ESR-1200};
\draw[link] (internet) -- node[right, font=\scriptsize] {WAN} (router);

% Уровень ядра (Серверная)
\node[core, below=of router, yshift=-0.5cm] (core_sw) {Ядро: MES5310-48};
\draw[link] (router) -- node[right, font=\scriptsize] {VLAN 10, 20} (core_sw);

% Серверы
\node[server, left=of core_sw, xshift=-2.5cm] (srv_cicd) {CI/CD С2122АА};
\node[server, above=of srv_cicd, yshift=-0.5cm] (srv_media) {Медиасервер С2242И};
\node[server, below=of srv_cicd, yshift=0.5cm] (srv_nas) {NAS QSRV-260202};

\draw[link] (srv_media.east) -- ++(1,0) |- node[pos=0.25, above, font=\scriptsize] {VLAN 40, 50} (core_sw.west);
\draw[link] (srv_cicd.east) -- node[above, font=\scriptsize] {VLAN 50} (core_sw.west);
\draw[link] (srv_nas.east) -- ++(1,0) |- node[pos=0.25, below, font=\scriptsize] {VLAN 50} (core_sw.west);

% Рамка для серверной
\draw[dashed, draw=gray, rounded corners] 
    ([xshift=-0.5cm,yshift=0.5cm]srv_media.north west |- router.north) rectangle 
    ([xshift=0.5cm,yshift=-0.5cm]core_sw.south east |- srv_nas.south) 
    node[above right, text=gray] at ([xshift=-0.5cm,yshift=0.5cm]srv_media.north west |- router.north) {\textbf{Серверная (пом. 19)}};

% Уровень доступа (Кроссовая)
\node[acc, below=of core_sw, yshift=-3cm, xshift=-2cm] (acc_poe) {
    \begin{tabular}{c}
        \textbf{Доступ (PoE+)} \\
        MES2448P
    \end{tabular}
};
\node[acc, left=of acc_poe, xshift=-0.5cm] (acc_poe2) {
    \begin{tabular}{c}
        \textbf{Доступ (PoE++)} \\
        AQ-N3000
    \end{tabular}
};
\node[acc, right=of acc_poe, xshift=0.5cm] (acc_base1) {
    \begin{tabular}{c}
        \textbf{Доступ (Б/П) №1} \\
        MES2448B
    \end{tabular}
};
\node[acc, right=of acc_base1, xshift=0.5cm] (acc_base2) {
    \begin{tabular}{c}
        \textbf{Доступ (Б/П) №2} \\
        MES2448B
    \end{tabular}
};

% Стекирование (кольцо)
\draw[thick, orange, <->] (acc_poe2.east) -- node[above, font=\scriptsize, text=orange] {Стек} (acc_poe.west);
\draw[thick, orange, <->] (acc_poe.east) -- node[above, font=\scriptsize, text=orange] {Стек} (acc_base1.west);
\draw[thick, orange, <->] (acc_base1.east) -- node[above, font=\scriptsize, text=orange] {Стек} (acc_base2.west);
\draw[thick, orange, <->] (acc_base2.south) -- ++(0,-0.4) -| node[pos=0.25, below, font=\scriptsize, text=orange] {Стек (Ring)} (acc_poe2.south);

% Рамка для логического стека
\draw[draw=orange!80, thick, rounded corners] 
    ([xshift=-0.3cm,yshift=0.6cm]acc_poe2.north west) rectangle 
    ([xshift=0.3cm,yshift=-0.9cm]acc_base2.south east)
    node[above right, text=orange!80] at ([xshift=-0.3cm,yshift=0.6cm]acc_poe2.north west) {\textbf{Стек коммутаторов доступа}};

% Связи доступа с ядром (LACP)
\path (acc_poe.north) -- (acc_base1.north) coordinate[midway] (stack_top);
\draw[uplink, ultra thick, blue] ([yshift=0.6cm]stack_top) -- node[right, font=\scriptsize, align=left] {LACP $4 \times 10$G SFP+ \\ VLAN 10, 20, 30, 40} (core_sw.south);

% Рамка для кроссовой
\draw[dashed, draw=gray, rounded corners] 
    ([xshift=-0.8cm,yshift=1.5cm]acc_poe2.north west) rectangle 
    ([xshift=0.8cm,yshift=-1.4cm]acc_base2.south east)
    node[above right, text=gray] at ([xshift=-0.8cm,yshift=1.5cm]acc_poe2.north west) {\textbf{Кроссовая (пом. 16)}};

\end{tikzpicture}
}
\end{figure}

\vspace{1cm}
\noindent Оптические аплинки 10G от стека коммутаторов доступа подключаются к портам SFP+ коммутатора ядра, формируя единый агрегированный логический канал (LACP / IEEE 802.3ad) емкостью 40 Гбит/с. Стекирование уровня доступа повышает отказоустойчивость (при выходе из строя одного коммутатора или линка стек продолжает функционировать) и упрощает администрирование всей группы портов.

\end{appendices}
\end{document}
